\documentclass[12pt,letterpaper]{article}
\usepackage{graphicx,textcomp}
\usepackage{natbib}
\usepackage{setspace}
\usepackage{fullpage}
\usepackage{color}
\usepackage[reqno]{amsmath}
\usepackage{amsthm}
\usepackage{fancyvrb}
\usepackage{amssymb,enumerate}
\usepackage[all]{xy}
\usepackage{endnotes}
\usepackage{lscape}
\newtheorem{com}{Comment}
\usepackage{float}
\usepackage{hyperref}
\newtheorem{lem} {Lemma}
\newtheorem{prop}{Proposition}
\newtheorem{thm}{Theorem}
\newtheorem{defn}{Definition}
\newtheorem{cor}{Corollary}
\newtheorem{obs}{Observation}
\usepackage[compact]{titlesec}
\usepackage{dcolumn}
\usepackage{tikz}
\usetikzlibrary{arrows}
\usepackage{multirow}
\usepackage{xcolor}
\newcolumntype{.}{D{.}{.}{-1}}
\newcolumntype{d}[1]{D{.}{.}{#1}}
\definecolor{light-gray}{gray}{0.65}
\usepackage{url}
\usepackage{listings}
\usepackage{color}

\definecolor{codegreen}{rgb}{0,0.6,0}
\definecolor{codegray}{rgb}{0.5,0.5,0.5}
\definecolor{codepurple}{rgb}{0.58,0,0.82}
\definecolor{backcolour}{rgb}{0.95,0.95,0.92}

\lstdefinestyle{mystyle}{
	backgroundcolor=\color{backcolour},   
	commentstyle=\color{codegreen},
	keywordstyle=\color{magenta},
	numberstyle=\tiny\color{codegray},
	stringstyle=\color{codepurple},
	basicstyle=\footnotesize,
	breakatwhitespace=false,         
	breaklines=true,                 
	captionpos=b,                    
	keepspaces=true,                 
	numbers=left,                    
	numbersep=5pt,                  
	showspaces=false,                
	showstringspaces=false,
	showtabs=false,                  
	tabsize=2
}
\lstset{style=mystyle}
\newcommand{\Sref}[1]{Section~\ref{#1}}
\newtheorem{hyp}{Hypothesis}

\title{Problem Set 2}
\date{Due: October 23, 2025}
\author{Lisanne van Vucht}

\begin{document}
	\maketitle
	\section*{Instructions}
\begin{itemize}
	\item Please show your work! You may lose points by simply writing in the answer. If the problem requires you to execute commands in \texttt{R}, please include the code you used to get your answers. Please also include the \texttt{.R} file that contains your code. If you are not sure if work needs to be shown for a particular problem, please ask.
	\item Your homework should be submitted electronically on GitHub.
	\item This problem set is due before 23:59 on Thursday October 23, 2025. No late assignments will be accepted.

\end{itemize}

	
	\vspace{.5cm}
	\section*{Question 1: Political Science}
		\vspace{.25cm}
	The following table was created using the data from a study run in a major Latin American city.\footnote{Fried, Lagunes, and Venkataramani (2010). ``Corruption and Inequality at the Crossroad: A Multimethod Study of Bribery and Discrimination in Latin America. \textit{Latin American Research Review}. 45 (1): 76-97.} As part of the experimental treatment in the study, one employee of the research team was chosen to make illegal left turns across traffic to draw the attention of the police officers on shift. Two employee drivers were upper class, two were lower class drivers, and the identity of the driver was randomly assigned per encounter. The researchers were interested in whether officers were more or less likely to solicit a bribe from drivers depending on their class (officers use phrases like, ``We can solve this the easy way'' to draw a bribe). The table below shows the resulting data.

\newpage
\begin{table}[h!]
	\centering
	\begin{tabular}{l | c c c }
		& Not Stopped & Bribe requested & Stopped/given warning \\
		\\[-1.8ex] 
		\hline \\[-1.8ex]
		Upper class & 14 & 6 & 7 \\
		Lower class & 7 & 7 & 1 \\
		\hline
	\end{tabular}
\end{table}

\begin{enumerate}
	
	\item [(a)]
	Calculate the $\chi^2$ test statistic by hand/manually (even better if you can do "by hand" in \texttt{R}).\\
	
\textit{First,	I created a matrix with the observed values. Second, I created a matrix with the expected values of each cell, by multiplying the sum of the row with the sum of the column and divide it by the total amount (grand total)of observations. Third, I compared the observed values with the expected to calculate $\chi^2$.  Fourth, I have calculated the expected values for each cell,  To calculate $\chi^2$ I used:
	 $\chi^2 = \sum{\frac{(f_{observed}-f_{expected})^2}{f_{expected}}}$}
	\begin{Verbatim}
		#1.  
		data <- matrix(NA, nrow= 2, ncol = 3)
		rownames(data) <- c("upper class", "lower class")
		colnames(data) <- c("not stopped", "bribe requested",
		"stopped/given warning")
		data[1,] <- c(14,6,7)
		data[2,] <- c(7,7,1)
		
		#2
		Rsum <-rowSums(data)
		Csum <- colSums(data)
		grand_total <-sum(data)
		expected <- outer(Rsum, Csum)/grand_total
		
		#3
		chi_sq <- sum((data - expected)^2/expected)
		chi_sq
	\end{Verbatim}
		
\textit{Result:  $\chi^2$ value is 3.791168. This is a relatively small test statistic meaning the average difference between the observed and expected values is modest. Now I will calculate the p-value to make a conclusion if alpha would be 0.1}
	'
	\item [(b)]
	Now calculate the p-value from the test statistic you just created (in \texttt{R}).\footnote{Remember frequency should be $>$ 5 for all cells, but let's calculate the p-value here anyway.}  What do you conclude if $\alpha = 0.1$?\\
	
	\begin{Verbatim}
	p_value <- pchisq(chi_sq,2, lower.tail = FALSE)
	p_value
\end{Verbatim}


\textit{The p-Value = 0,1502. This is larger than our alpha 0.1. Therefore we fail to reject H0. Meaning we do not have sufficient evidence to conclude a statistically independent difference between upper and lower class drivers.}
	
	\item [(c)] Calculate the standardized residuals for each cell and put them in the table below.
	\vspace{1cm}
	
	\begin{table}[h]
		\centering
		\begin{tabular}{l | c c c }
			& Not Stopped & Bribe requested & Stopped/given warning \\
			\\[-1.8ex] 
			\hline \\[-1.8ex]
			Upper class  & 0.1360828   & 0.1360828   & 0.1360828  \\
			\\
			Lower class & -0.1825742   &   1.0939393  &   -1.098701 \\
			
		\end{tabular}
	\end{table}
	
\textit{	I have calculated the residuals using the formula: $z =\frac{f_{observed}-f_{expected}}{se}$}

	\item [(d)] How might the standardized residuals help you interpret the results?  
	
	\textit{With standardized residuals ranging from –1.10 to +1.09, all of the observed values are close to their expected values if H0 would be true. Meaning, none of the  combinations of class (upper vs. lower) and outcome (not stopped, bribe requested, stopped/given warning) stand out in as regard to having unusually large deviations. 
		This alligns with the previous conclusion regarding the p-value. }
	
\end{enumerate}

\section*{Question 2: Economics}
Chattopadhyay and Duflo were interested in whether women promote different policies than men.\footnote{Chattopadhyay and Duflo. (2004). ``Women as Policy Makers: Evidence from a Randomized Policy Experiment in India. \textit{Econometrica}. 72 (5), 1409-1443.} Answering this question with observational data is pretty difficult due to potential confounding problems (e.g. the districts that choose female politicians are likely to systematically differ in other aspects too). Hence, they exploit a randomized policy experiment in India, where since the mid-1990s, $\frac{1}{3}$ of village council heads have been randomly reserved for women. A subset of the data from West Bengal can be found at the following link: \url{https://raw.githubusercontent.com/kosukeimai/qss/master/PREDICTION/women.csv}\\

\noindent Each observation in the data set represents a village and there are two villages associated with one GP (i.e. a level of government is called "GP"). Figure~\ref{fig:women_desc} below shows the names and descriptions of the variables in the dataset. The authors hypothesize that female politicians are more likely to support policies female voters want. Researchers found that more women complain about the quality of drinking water than men. You need to estimate the effect of the reservation policy on the number of new or repaired drinking water facilities in the villages.
\vspace{.5cm}
\begin{figure}[h!]
	\caption{\footnotesize{Names and description of variables from Chattopadhyay and Duflo (2004).}}
	\vspace{.5cm}
	\centering
	\label{fig:women_desc}
\end{figure}		

\newpage
\begin{enumerate}
	\item [(a)] State a null and alternative (two-tailed) hypothesis. 
	

	H0: There is no difference in the mean number of drinking water facilities between villages with and without a reserved female GP.
	
	Ha: There is a difference in the mean number of drinking water facilities between villages with and without a reserved female GP.
	
	In other words, under H0 the average number of facilities is the same in reserved and non-reserved villages. If the treatment "reserved" had no effect, the mean of water facilities for "reserved = 0" and "reserved = 1" would be identical.
	
	Our null hypothesis states that there is no difference in the mean number of new or repaired drinking water facilities between villages that received treatment ("reserved") and those that did not.
	
	
	\vspace{6cm}
	\item [(b)] Run a bivariate regression to test this hypothesis in \texttt{R} (include your code!).
	
	\begin{Verbatim}
		download.file(
		url = "https://raw.githubusercontent.com
		/kosukeimai/qss/master/PREDICTION/women.csv",
		destfile = "women.csv",
		mode = "wb"
		)
		
		women
		
		women <- read.csv("women.csv")
		
		#check structure of data before modelling 
		
		str(women)
		head(women)
		summary(women)
		
		model <- lm(water ~ reserved, data = women)
		summary(model)

		Results:
		
		Residuals:
		Min      1Q  Median      3Q     Max 
		-23.991 -14.738  -7.865   2.262 316.009 
		
		Coefficients:
		Estimate Std. Error t value Pr(>|t|)    
		(Intercept)   14.738      2.286   6.446 4.22e-10 ***
		reserved       9.252      3.948   2.344   0.0197 *  
		---
		Signif. codes:  0 ‘***’ 0.001 ‘**’ 0.01 ‘*’ 0.05 ‘.’ 0.1 ‘ ’ 1
		
		Residual standard error: 33.45 on 320 degrees of freedom
		Multiple R-squared:  0.01688,	Adjusted R-squared:  0.0138 
		F-statistic: 5.493 on 1 and 320 DF,  p-value: 0.0197
		
				\end{Verbatim}
	
	\vspace{6cm}
	\item [(c)] Interpret the coefficient estimate for reservation policy. 
	
	The coefficient estimate for the reservation policy is 9.252. This means that on average, we can expect villages with reserved female council heads to have 23.9 more new or repaired drinking water facilities. 

	Although this effect is statistically significant, p = 0.0197, the residual standard error of 33.45 shows that  on avarage there’s a lot of variation in the number of facilities that isn’t explained by the treatment alone. The Rsquared of 0.017 confirms that only a small portion of the total variance is accounted for by the reservation policy.
	
	In sum, therefore looking at the avarage number of drinking water facilities and reservation policy the effect is modest. This means the reservation policy has a statistically significant positive impact on avarage, but  other factors may also influence the number of water facilities in villages.
	
	
\end{enumerate}



%\end{enumerate}
\end{document}
